
\documentclass[12pt]{article}
%%%%%%%%%%%%%%%%%%%%%%%%%%%%%%%%%%%%%%%%%%%%%%%%%%%%%%%%%%%%%%%%%%%%%%%%%%%%%%%%%%%%%%%%%%%%%%%%%%%%%%%%%%%%%%%%%%%%%%%%%%%%%%%%%%%%%%%%%%%%%%%%%%%%%%%%%%%%%%%%%%%%%%%%%%%%%%%%%%%%%%%%%%%%%%%%%%%%%%%%%%%%%%%%%%%%%%%%%%%%%%%%%%%%%%%%%%%%%%%%%%%%%%%%%%%%
\usepackage{amsmath}
\usepackage[compat2]{geometry}

\setcounter{MaxMatrixCols}{10}
%TCIDATA{TCIstyle=LaTeX article (bright).cst}

%TCIDATA{OutputFilter=LATEX.DLL}
%TCIDATA{Version=5.50.0.2953}
%TCIDATA{<META NAME="SaveForMode" CONTENT="1">}
%TCIDATA{BibliographyScheme=Manual}
%TCIDATA{Created=Tuesday, August 14, 2012 13:45:33}
%TCIDATA{LastRevised=Friday, April 08, 2022 10:30:56}
%TCIDATA{<META NAME="GraphicsSave" CONTENT="32">}
%TCIDATA{<META NAME="DocumentShell" CONTENT="Exams and Syllabi\SW\Assignment">}

\setlength{\topmargin}{-1.0in}
\setlength{\textheight}{9.25in}
\setlength{\oddsidemargin}{0.0in}
\setlength{\evensidemargin}{0.0in}
\setlength{\textwidth}{6.5in}
\def\labelenumi{\arabic{enumi}.}
\def\theenumi{\arabic{enumi}}
\def\labelenumii{(\alph{enumii})}
\def\theenumii{\alph{enumii}}
\def\p@enumii{\theenumi.}
\def\labelenumiii{\arabic{enumiii}.}
\def\theenumiii{\arabic{enumiii}}
\def\p@enumiii{(\theenumi)(\theenumii)}
\def\labelenumiv{\arabic{enumiv}.}
\def\theenumiv{\arabic{enumiv}}
\def\p@enumiv{\p@enumiii.\theenumiii}
\pagestyle{plain}
\setcounter{secnumdepth}{0}
\parindent=0pt
\input{tcilatex}
\begin{document}


\begin{center}
\begin{equation*}
\FRAME{itbpF}{1.3214in}{0.9496in}{0in}{}{}{Figure}{\special{language
"Scientific Word";type "GRAPHIC";maintain-aspect-ratio TRUE;display
"USEDEF";valid_file "T";width 1.3214in;height 0.9496in;depth
0in;original-width 1.3552in;original-height 0.9669in;cropleft "0";croptop
"1";cropright "1";cropbottom "0";tempfilename
'QYQ9VJ00.bmp';tempfile-properties "XPR";}}
\end{equation*}%
\textbf{Electromagnetismo - 230045}

\emph{Gu\'{\i}a N}$%
%TCIMACRO{\U{b0}}%
%BeginExpansion
{{}^\circ}%
%EndExpansion
1$\emph{\ - Fuerza el\'{e}ctrica entre cargas puntuales}

\textit{Profesor: Arturo Fern\'{a}ndez P\'{e}rez}
\end{center}

\section{Problemas unidimensionales}

\begin{enumerate}
\item Una carga puntual $\ q_{1}=$ 85~$[nC]$ se ubica en el punto $x=-1$ [m]
y otra carga $q_{2}=-84$ [$\mu C$] se ubica en el punto $x=0.85$ [m].\
Determine la fuerza el\'{e}ctrica ejercida sobre la carga $q_{2}.$ \textbf{R:%
} $\vec{F}=0,0188$ $(-\hat{\imath})$ [N]

\item Una varilla de vidrio se somete a frotaci\'{o}n y adquiere una carga
neta $q_{1}=-6$ $[mC].$ Se acerca a una varilla pl\'{a}stica que tiene una
carga $q_{2}=85$ [$\mu C$]. Determine la distancia a la cu\'{a}l deben estar
para que la fuerza el\'{e}ctrica entre ellas tenga magnitud $5.1\times
10^{-3}$ [N]. \textbf{R:} $948.68$ [m]

\item Dado el sistema de la Fig. 1, determine la fuerza el\'{e}ctrica que
siente una carga $q=5$ $[nC]$, si \'{e}sta se ubica en el punto $A$ y luego
en el punto $B$. \textbf{R:} $A:$ $\vec{F}=4,375\times 10^{-5}(\hat{\imath})$
[N] \ ; $B:\vec{F}=3,272\times 10^{-5}(\hat{\imath})$ [N] 
\begin{equation*}
\FRAME{itbpFU}{2.2632in}{0.6806in}{0in}{\Qcb{Figura 1}}{}{Figure}{\special%
{language "Scientific Word";type "GRAPHIC";maintain-aspect-ratio
TRUE;display "USEDEF";valid_file "T";width 2.2632in;height 0.6806in;depth
0in;original-width 4.8205in;original-height 1.4313in;cropleft "0";croptop
"1";cropright "1";cropbottom "0";tempfilename
'QYQ9VJ01.wmf';tempfile-properties "XPR";}}
\end{equation*}

\item Dos cargas puntuales, como muestra la Fig. 2, est\'{a}n ubicadas sobre
el eje X. Encuentre el punto sobre este eje donde al colocar una carga
puntual de magnitud $q=3$ $\mu C$ la fuerza neta sobre ella sea nula. 
\textbf{R:} A $3.275$ [m] a la izquierda de la carga $+2,5\mu C$ 
\begin{equation*}
\FRAME{itbpFU}{2.0868in}{0.5604in}{0in}{\Qcb{Figura 2}}{}{Figure}{\special%
{language "Scientific Word";type "GRAPHIC";maintain-aspect-ratio
TRUE;display "USEDEF";valid_file "T";width 2.0868in;height 0.5604in;depth
0in;original-width 4.209in;original-height 1.1113in;cropleft "0";croptop
"1";cropright "1";cropbottom "0";tempfilename
'QYQ9VJ02.wmf';tempfile-properties "XPR";}}
\end{equation*}%
\pagebreak
\end{enumerate}

\section{Problemas bidimensionales}

\begin{enumerate}
\item En los v\'{e}rtices de un cuadrado de lado $a=5.5$ [cm] se ubican
cuatro cargas puntuales id\'{e}nticas $q=-5$ [$\mu C].$ Determine la fuerza
el\'{e}ctrica inducida sobre la carga ubicada en la esquina inferior
izquierda del cuadrado. \textbf{R: }$\vec{F}=100,678\left( -\hat{\imath}-%
\hat{\jmath}\right) $ [N]

\item Dos cargas puntuales de igual carga $(\left\vert q_{1}\right\vert
=\left\vert q_{2}\right\vert =5$ $[nC])$ pero signo contrario (un \textit{%
dipolo el\'{e}ctrico}) se ubican como muestra la Fig. 3. Determine la fuerza
el\'{e}ctrica que se induce sobre una carga puntual $q_{0}=3$ $[nC]\,,$ si
ella se ubica en los puntos $a$, $b$, y $c$. \textbf{R: }$\vec{F}%
_{a}=1,219\times 10^{-4}\hat{\imath}$ [N]$;\vec{F}_{b}=7,749\times
10^{-5}\left( -\hat{\imath}\right) $ [N] ; $\vec{F}_{c}=6,144\times 10^{-6}%
\hat{\imath}$ [N] 
\begin{equation*}
\FRAME{itbpFU}{2.2303in}{2.079in}{0in}{\Qcb{Figura 3}}{}{Figure}{\special%
{language "Scientific Word";type "GRAPHIC";maintain-aspect-ratio
TRUE;display "USEDEF";valid_file "T";width 2.2303in;height 2.079in;depth
0in;original-width 5.1534in;original-height 4.8066in;cropleft "0";croptop
"1";cropright "1";cropbottom "0";tempfilename
'QYQ9VJ03.wmf';tempfile-properties "XPR";}}
\end{equation*}

\item Si dos electrones est\'{a}n a una distancia de $1.5\times 10^{-10}$
[m] ($1.5$ $\mathring{A}$) de un\ prot\'{o}n, como muestra la Fig. 4,
determine la fuerza neta ejercida sobre el prot\'{o}n. \textbf{R: }$\vec{F}%
=1,457\times 10^{-8}\hat{\imath}+9,281\times 10^{-9}\hat{\jmath}$ [N]%
\begin{equation*}
\FRAME{itbpFU}{1.0741in}{0.9366in}{0in}{\Qcb{Figura 4}}{}{Figure}{\special%
{language "Scientific Word";type "GRAPHIC";maintain-aspect-ratio
TRUE;display "USEDEF";valid_file "T";width 1.0741in;height 0.9366in;depth
0in;original-width 2.0427in;original-height 1.7789in;cropleft "0";croptop
"1";cropright "1";cropbottom "0";tempfilename
'QYQ9VJ04.wmf';tempfile-properties "XPR";}}
\end{equation*}

\item Dos cargas puntuales $q_{1}$ y $q_{2}$ est\'{a}n ubicadas a 4.5 $[cm]$
de distancia. Otra carga puntual $Q=-1.75$ $\mu C$ y masa $m=5$ [g] est\'{a}
ubicada a 3 $[cm]$ de c/u de las cargas y se suelta desde el reposo, como
muestra la Fig. 5. Se observa que ella adquiere una aceleraci\'{o}n de 324 $%
[m/s^{2}]\,\ $hacia arriba, en una direcci\'{o}n paralela a la l\'{\i}nea
que conecta las dos cargas puntuales. Determine $q_{1}$ y $q_{2}.$ \textbf{%
R: }$\left\vert q_{1}\right\vert =\left\vert q_{2}\right\vert =61,74$ $\mu C$%
\begin{equation*}
\FRAME{itbpFU}{1.1268in}{1.529in}{0in}{\Qcb{Figura 5}}{}{Figure}{\special%
{language "Scientific Word";type "GRAPHIC";maintain-aspect-ratio
TRUE;display "USEDEF";valid_file "T";width 1.1268in;height 1.529in;depth
0in;original-width 2.1525in;original-height 2.9317in;cropleft "0";croptop
"1";cropright "1";cropbottom "0";tempfilename
'QYQ9VJ05.wmf';tempfile-properties "XPR";}}
\end{equation*}

\item Tres cargas puntuales est\'{a}n en los v\'{e}rtices de un tri\'{a}%
ngulo is\'{o}sceles, como muestra la Fig. 6. Las cargas $\pm 5$ $\mu C$
forman un dipolo el\'{e}ctrico. Encuentre la fuerza que ejerce la carga de
-10 $\mu C$ sobre el dipolo. \textbf{R:} $\vec{F}=1687,49\hat{\jmath}$ [N] 
\begin{equation*}
\FRAME{itbpFU}{1.5273in}{1.5766in}{0in}{\Qcb{Figura 6}}{}{Figure}{\special%
{language "Scientific Word";type "GRAPHIC";maintain-aspect-ratio
TRUE;display "USEDEF";valid_file "T";width 1.5273in;height 1.5766in;depth
0in;original-width 2.9871in;original-height 3.0839in;cropleft "0";croptop
"1";cropright "1";cropbottom "0";tempfilename
'QYQ9VJ06.wmf';tempfile-properties "XPR";}}
\end{equation*}
\end{enumerate}

\end{document}
